\scanpage{132}\mspage{- 116 -}\hfill IX

\uline{dans la catégorie des extensions panachées de} $E_\eta$ \uline{par} $F_\eta$. \uline{Pour tout groupe de \unsure{Barsotti-Tate|Barootti-Tate} extension panachée} $X_\eta$ \uline{de} $E_\eta$ \uline{par} $F_\eta$, \uline{on considère l'accouplement correspondant}
\[
 Q_{\unsure{o|0}}\otimes Q_o^*\longrightarrow\mathbb Z_{\ell s}\ , \tag*{(9.5.6.2)}
\]
\uline{induit par l'accouplement du type} (9.5.5) \uline{défini sur le} hensélisé \uline{de} $S$. \uline{Ceci posé}, \uline{le foncteur} (9.5.6.1) \uline{est pleinement fidèle}, \uline{et son image essentielle est formée des groupes de \deleted{\illegible{1}}Barsotti-Tate} $X_\eta$ \uline{extensions panachées de} $E_\eta$ \uline{par} $F_\eta$ \uline{telles que l'accouplement} (9.5.6.2) \uline{correspondant soit nul}. \uline{De plus}, \uline{lorsque} $\mathrm{car}.K=0$ (\uline{de sorte qu'on peut appliquer le théorème de TATE} [33 th. 4]), \uline{alors} $X_\eta$ \uline{appartient à l'image essentielle si et seulement si} $X_\eta$ \uline{se prolonge en un groupe de Barsotti-Tate sur} $S$.

9.5.6.3.\quad Nous laissons au lecteur le soin d'étendre cet énoncé au cas où on suppose seulement que $S$ est un schéma noethérien régulier connexe de dimension 1.

La première assertion de 9.5.6 est une conséquence immédiate de 9.4.6, appliqué aux composants des groupes de Barsotti-Tate envisagés. Il reste à prouver que si un groupe de \deleted{\illegible{1}}Barsotti-Tate $X_\eta$, extension panachée de $E_\eta$ par $F_\eta$, se prolonge en un groupe de Barsotti-Tate $X_\eta$ sur $S$, il se prolonge aussi comme extension panachée. Or en vertu du théorème de TATE [33], l'homomorphisme $X_\eta\longrightarrow Q_\eta$ se prolonge de façon unique en un homomorphisme
\[
 X\longrightarrow Q\quad .
\]
% archive sheet number: 131
% stamp: numbered label 434.
